%%%%%%%%%%%%%%%%%%%%%%%%%%%%%%%%%%%%%%%%%%%%%%%%%%%%%%%%%%%%%%%%%%%%%%
% Lớp 12 | Ứng dụng đạo hàm để khảo sát hàm số | Khảo sát sự biến thiên và vẽ đồ thị hàm số
% Dạng 2D1?5-2: Các phép biến đổi đồ thị
% Số câu: 2  (sinh tự động bởi tach_cau.py; từ nay thầy sửa trực tiếp file này)
%%%%%%%%%%%%%%%%%%%%%%%%%%%%%%%%%%%%%%%%%%%%%%%%%%%%%%%%%%%%%%%%%%%%%%

%%%%%%%%%%%%%%%% TRẮC NGHIỆM 4 PHƯƠNG ÁN %%%%%%%%%%%%%%%%

% nguồn: Phan1_Lop_12_Dai_so_Chuong_1
\begin{ex}%[2D1H5-2]%[ma:c8906312]
		Cho hàm số $y=f(x)$ xác định, liên  tục trên $\mathbb{R}$ và có đồ thị như hình bên.
		\begin{center}
			\begin{tikzpicture}[line join=round,line cap=round,>=stealth]
				\tikzset{ethicknode/.style={scale=0.9}} 
				\draw[->] (-2.5,0)--(2,0) node[below] {$x$}; 
				 \draw[->] (0,-1.1)--(0,3.1) node[below left] {$y$};
				\draw (0,0) node [below left] {$O$}; 
				\foreach \x/\nx in  {1/1}
				\draw[thin] (\x,1pt)--(\x,-1pt) node  [below] {$\nx$};
				\foreach \y/\ny  in {1/1}
				\draw[thin] (1pt,\y)--(-1pt,\y) node  [left] {$\ny$};
				\begin{scope}
					\clip (-2.5,-1) rectangle (1.5,3); 
					 \draw[samples=200,domain=-2.5:2.5,smooth,variable=\x] plot (\x,{3^(\x)});
				\end{scope}
			\end{tikzpicture}
		\end{center}
		Đồ thị nào dưới đây là đồ thị của  hàm số $y=f(x)+1$?
	\choice 
			{\begin{tikzpicture}[line join=round,line cap=round,>=stealth]
		 \tikzset{ethicknode/.style={scale=0.9}}
		\draw[->] (-2.5,0)--(2,0) node[below] {$x$}; 
		\draw[->]  (0,-1.1)--(0,3.1) node[below left] {$y$};
		\draw (0,0) node [below left] {$O$}; 
		 \foreach \x/\nx in {1/1}
		\draw[thin] (\x,1pt)--(\x,-1pt) node [below] {$\nx$}; 
		\foreach \y/\ny in {1/1} 
		 \draw[thin] (1pt,\y)--(-1pt,\y) node [left] {$\ny$};
		\begin{scope}
			\clip  (-2.5,-1) rectangle (1.5,3);
			\draw[samples=200,domain=-2.5:2.5,smooth,variable=\x] plot (\x,{3^(\x)}); 
		\end{scope}
	\end{tikzpicture}
			} 
			{\True \begin{tikzpicture}[line join=round,line cap=round,>=stealth]
		 \tikzset{ethicknode/.style={scale=0.9}}
		 \draw[->] (-2.5,0)--(2,0) node[below] {$x$};
		\draw[->] (0,-1.1)--(0,3.1)  node[below left] {$y$};
		\draw (0,0) node [below left] {$O$}; 
		 \foreach \x/\nx in {1/1}
		\draw[thin] (\x,1pt)--(\x,-1pt)  node [below] {$\nx$};
		\foreach \y/\ny in  {1/1,2/2}
		\draw[thin] (1pt,\y)--(-1pt,\y) node  [left] {$\ny$};
		\begin{scope}
			\clip  (-2.5,-1) rectangle (1.5,3);
			\draw[samples=200,domain=-2.5:2.5,smooth,variable=\x]  plot (\x,{3^(\x)+1});
		\end{scope}
	\end{tikzpicture}
			 }
			{\begin{tikzpicture}[line join=round,line cap=round,>=stealth]
		 \tikzset{ethicknode/.style={scale=0.9}}
		 \draw[->] (-2.5,0)--(2,0) node[below] {$x$};
		 \draw[->] (0,-1.1)--(0,3.1) node[below left] {$y$};
		\draw (0,0) node  [below right] {$O$};
		\foreach  \x/\nx in {1/1}
		\draw[thin] (\x,1pt)--(\x,-1pt)  node [below] {$\nx$};
		\foreach  \y/\ny in {1/1}
		\draw[thin] (1pt,\y)--(-1pt,\y) node [left]  {$\ny$};
		\begin{scope}
			\clip (-2.5,-1) rectangle  (1.5,3);
			\draw[samples=200,domain=-2.5:2.5,smooth,variable=\x] plot (\x,{3^(\x)-1}); 
		\end{scope}
	\end{tikzpicture}
			 }
			{\begin{tikzpicture}[line join=round,line cap=round,>=stealth]
		\tikzset{ethicknode/.style={scale=0.9}} 
		\draw[->]  (-2.5,0)--(2,0) node[below] {$x$};
		\draw[->] (0,-1.1)--(0,3.1) node[below left]  {$y$};
		\draw (0,0) node  [below left] {$O$};
		 \foreach \x/\nx in {1/1}
		\draw[thin] (\x,1pt)--(\x,-1pt)  node [below] {$\nx$};
		\foreach \y/\ny  in {1/1,2/2}
		 \draw[thin] (1pt,\y)--(-1pt,\y) node [left] {$\ny$};
		\begin{scope}
			 \clip (-2.5,-1) rectangle (1.5,3);
			 \draw[samples=200,domain=-2.5:2.5,smooth,variable=\x] plot (\x,{3^(\x+0.7)});
		\end{scope}
	\end{tikzpicture}
			} 
	\loigiai{ 
			Từ đồ thị hàm  số $y=f(x)$ tìm đồ thị hàm số $y=f(x)+a$.\\
			Tịnh tiến đồ thị hàm  số $y=f(x)$ lên phía trên (theo phương $Oy$) $a$ đơn vị nếu $a>0$.\\
			Hàm số $y=f(x)+1 \Rightarrow a=1>0$ nên ta được đồ thị hàm số cần tìm tịnh tiến đồ  thị hàm số $y=f(x)$ lên phía trên (theo phương $Oy$) $1$ đơn vị.
		} 
	\end{ex}


%%%%%%%%%%%%%%%% ĐÚNG/SAI %%%%%%%%%%%%%%%%

% nguồn: Phan1_Lop_12_Dai_so_Chuong_1
\begin{ex}%[2D1H5-2]%[ma:a6eee187]
	Cho hàm số $y=\dfrac{x^2+4x+5}{x+2}$. Khi đó 
	 \choiceTF
		{\True Hàm  số có tập xác định $\mathscr{D}=\mathbb{R}\setminus\{-2\}$}
		{Hàm  số đồng biến trên khoảng $(-\infty;-2)$ và $(-2; 3)$}
		 {Hàm số có tiệm cận xiên là $y=x+3$}
		{\True Đồ thị hàm số có tâm  đối xứng nằm trên đường thẳng $2x-y+4=0$}
	 \loigiai{
		\begin{itemchoice}
			 \itemch Đúng. Tập xác định $\mathscr{D}=\mathbb{R} \setminus\{-2\}$.
			\itemch Sai. Ta có $y=\dfrac{x^2+4x+5}{x+2}=x+2+\dfrac{1}{x+2}$. Khi  đó $y'=\dfrac{x^2+4x+3}{(x+2)^2}$.\\
			Xét  $y'=0\Leftrightarrow\hoac{&x=-1\\
			 &x=-3.}$\\
			Ta có  bảng biến thiên
			\begin{center}
				\begin{tikzpicture}
					\tkzTabInit[nocadre=false,lgt=1.5,espcl=2.5,deltacl=0.6] 
						{$x$  /0.7,$f'(x)$ /0.7,$f (x)$ /2.}{$-\infty$,$-3$,$-2$,$-1$,$+\infty$}
					\tkzTabLine {,+,$0$,-,d,-,$0$,+,} 
					 \tkzTabVar {-/$-\infty$,+/$-2$,-D+/$-\infty$/$+\infty$,-/$2$,+/$+\infty$}
				\end{tikzpicture}
			\end{center}
			Dựa vào bảng biến  thiên ta có hàm số đồng biến trên khoảng $(-\infty;-3)$ và $(-1; 3)$.
			\itemch Sai. Ta có $y=\dfrac{x^2+4x+5}{x+2}=x+2+\dfrac{1}{x+2}$ nên $\lim\limits_{x\to+\infty}(y-x-2)=0$  và $\break \lim\limits_{x\to-\infty}(y-x-2)=0$ nên
			đồ  thị hàm số có tiệm cận xiên là $y=x+2$.
			\itemch Đúng. Ta thấy hàm số có tiệm cận tiệm  cận đứng là $x=-2$ và tiệm cận xiên là $y=x+2$ nên đồ thị hàm số có tâm đối xứng $I(-2; 0)$.\\
			Khi đó đồ thị hàm số có tâm đối xứng  $I(-2; 0)$ nằm trên đường thẳng $2x-y+4=0$.
		\end{itemchoice}
	 }
\end{ex}
